% Appendix for EpiContext paper
\section{Proof of Convergence}
\label{sec:proof}

\begin{proof}[Full proof of Theorem 1]
Consider a node $v_i$ with epigenetic weight $w_i^{(t)}$ at turn $t$. The update rule is:

\begin{equation}
    w_i^{(t+1)} = \text{clip}_{[0,1]}\left(w_i^{(t)} + \Delta w_i^{(t)}\right)
\end{equation}

where $\Delta w_i^{(t)} = +\delta_{\text{reward}}$ if $v_i$ was active and the turn succeeded, $-\delta_{\text{penalize}}$ if $v_i$ was active and the turn failed, and $+\delta_{\text{explore}}$ if $v_i$ was inactive.

Let $p_i = P(\text{turn succeeds} \mid v_i \text{ active})$ be the conditional success probability when $v_i$ is included in the context. The expected drift when $v_i$ is active is:

\begin{equation}
    \mathbb{E}[\Delta w_i \mid v_i \text{ active}] = p_i \delta_{\text{reward}} - (1-p_i) \delta_{\text{penalize}}
\end{equation}

Define the critical threshold $p^* = \frac{\delta_{\text{penalize}}}{\delta_{\text{reward}} + \delta_{\text{penalize}}}$. Then:

\begin{itemize}
    \item If $p_i > p^*$, then $\mathbb{E}[\Delta w_i] > 0$, and $w_i$ drifts upward toward 1.
    \item If $p_i < p^*$, then $\mathbb{E}[\Delta w_i] < 0$, and $w_i$ drifts downward toward 0.
    \item If $p_i = p^*$, the drift is zero and $w_i$ follows a symmetric random walk.
\end{itemize}

Since $w_i$ is bounded in $[0,1]$ and the step sizes $\delta_{\text{reward}}, \delta_{\text{penalize}}$ are constant, the process is a bounded random walk with reflecting boundaries at 0 and 1. By standard results for bounded random walks with constant drift, $w_i$ converges almost surely to 1 when $p_i > p^*$ and to 0 when $p_i < p^*$.

For inactive nodes, the exploration bonus $\delta_{\text{explore}}$ ensures that all nodes are periodically reactivated, preventing permanent silencing of potentially useful information. The frequency of reactivation is controlled by $\delta_{\text{explore}}$: larger values increase exploration at the cost of occasionally including low-quality nodes.

Under a stationary task distribution, the success probabilities $p_i$ are fixed, and the weights converge to a stationary distribution where $w_i \in \{0, 1\}$ for all $i$ with $p_i \neq p^*$. This establishes that \EpiContext{} learns to include exactly those context elements that contribute positively to task success.
\end{proof}

\section{Additional Experimental Results}

\subsection{Per-Turn Fitness Evolution}

Figure~\ref{fig:fitness_evolution} shows the evolution of fitness scores across turns for \EpiContext{} and baseline methods on a representative WebArena task. \EpiContext{} exhibits a characteristic learning curve: initial turns have moderate fitness as the system explores context configurations, followed by increasing fitness as successful patterns are reinforced.

% Prompt for Figure A1 (fitness evolution plot):
% A line chart showing fitness score (y-axis, range 0 to 1) vs turn number (x-axis, range 1 to 50).
% Four lines: EpiContext (purple, solid, upward trend), ReAct (orange, dashed, flat),
% MemGPT (green, dotted, slight upward), Full-Context (red, dash-dot, flat at low value).
% X-axis label: "Turn Number", Y-axis label: "Fitness Score F(P)".
% Legend in upper left. Professional, clean style suitable for academic publication.

\subsection{Computational Overhead Analysis}

Table~\ref{tab:overhead} reports the computational overhead of \EpiContext's auxiliary operations. The total overhead is less than 3\% of the token savings achieved.

\begin{table}[h]
\centering
\caption{Computational overhead breakdown (per 100 turns).}
\label{tab:overhead}
\begin{tabular}{lccc}
\toprule
\textbf{Operation} & \textbf{Calls/100 turns} & \textbf{Tokens/Call} & \textbf{Total Overhead} \\
\midrule
Methylation summarization & 5--8 & 200--500 & 1,000--4,000 \\
Acetylation relevance & 100 & 50--100 & 5,000--10,000 \\
Crossover evaluation & 2--5 & 500--1,000 & 1,000--5,000 \\
\midrule
\textbf{Total overhead} & & & \textbf{7,000--19,000} \\
\textbf{Token savings} & & & \textbf{300,000--500,000} \\
\bottomrule
\end{tabular}
\end{table}

\subsection{Implementation Details}

The complete \EpiContext{} implementation is available in the supplementary material. Key implementation details:

\begin{itemize}
    \item \textbf{Context Graph}: Implemented using Python dictionaries for $O(1)$ node lookup and adjacency lists for efficient traversal. Memory footprint scales as $O(|\mathcal{V}| + |\mathcal{E}|)$.
    \item \textbf{Epigenetic Tags}: Stored as a NumPy array of shape $(|\mathcal{V}|,)$ for vectorized batch updates. Tag updates use in-place operations for efficiency.
    \item \textbf{Serialization}: Full state serialization to JSON enables checkpointing and cross-session persistence of epigenetic marks.
    \item \textbf{Concurrency}: The crossover operator supports parallel variant evaluation through Python's concurrent.futures, enabling true parallel LLM calls in production deployments.
\end{itemize}